\section{Ternary incidence, discriminant gluing, and the Leech construction}
\label{sec:lf-foundations}

The lattice construction uses the incidence information that
accompanies the dodecaphasic record. Its domain preserves the ternary
words, the frame in which they are represented, the pre-carry
orientation, and the flag distinguishing one of the twelve positions. The
passage to a lattice transmits these data through successive
operations: linear encoding, incidence selection, discriminant
gluing, and formation of a neighbour. Lattice classification
enters after proving positivity, evenness, self-duality,
and the absence of roots in the resulting construction.

These operations continue the same APP--TRIT--TPK
generation and the same compatible refinement of the joint
continuum that precede numerical reading \cite{HMT-IV}.
The finite matrix, code, and metric defined below
are realizations with precise domains. The lattice
proof receives the marked state; the final scalar of a
constant plays the role of neither a glue code nor a length.

\subsection{Elliptic chart and incidence transport}
\label{subsec:lf-incidence}

The six units \(\{1,2,4,5,7,8\}\) modulo nine, preserved
by modular transport, admit a six-position chart.
We identify it with
\(\mathbb P^1(\mathbb F_5)=\{\infty,0,1,2,3,4\}\), in that order.
This identification fixes coordinates for realizing the quadratic
relation of the elliptic regime of TRIT.

Let \(\chi:\mathbb F_5\to\{0,1,-1\}\) be the quadratic character:
\(\chi(0)=0\), \(\chi(1)=\chi(4)=1\), and
\(\chi(2)=\chi(3)=-1\). We define the integer matrix
\(C_{\mathrm P}\) with zero diagonal,
\((C_{\mathrm P})_{\infty,a}=(C_{\mathrm P})_{a,\infty}=1\),
and finite entries \((C_{\mathrm P})_{a,b}=\chi(b-a)\).
Its square is \(5I_6\). Indeed, each diagonal entry of the
square sums five units. The off-diagonal entries involving
\(\infty\) vanish because the sum of
\(\chi\) is zero. For two distinct finite indices, the
contribution of one from position \(\infty\) cancels with
\(\sum_c\chi(c-a)\chi(c-b)=-1\). This last identity is verified
directly on the five residues, or through the affine
substitution taking \((a,b)\) to \((0,1)\).

Reduction modulo three produces the symmetric matrix
\begin{equation}
 A_W=
 \begin{pmatrix}
 0&1&1&1&1&1\\
 1&0&1&2&2&1\\
 1&1&0&1&2&2\\
 1&2&1&0&1&2\\
 1&2&2&1&0&1\\
 1&1&2&2&1&0
 \end{pmatrix},
 \qquad A_W^2=2I_6=-I_6
 \quad\text{in }\mathbb F_3.
 \label{eq:lf-aw}
\end{equation}
The equality preserves the elliptic type in this finite codomain.
The finite field and the subsequent Archimedean realization of the
imaginary unit have different domains.

For a matrix \(A\in M_6(\mathbb F_3)\) and a row word
\(w\in\mathbb F_3^6\), we write
\[
 \gamma_A(w)=(w,wA)\in\mathbb F_3^{12}.
\]
If \(f\in\operatorname{GL}_6(\mathbb F_3)\), then
\begin{equation}
 \gamma_{fAf^{-1}}(wf^{-1})
   =\gamma_A(w)(f^{-1}\oplus f^{-1}).
 \label{eq:lf-frame-covariance}
\end{equation}
The first half is \(wf^{-1}\); the second is
\(wf^{-1}fAf^{-1}=wAf^{-1}\), proving the identity.

Let \(X_n^{\mathrm{mar}}\) be the space of enriched prefixes of
length \(n\), with a declared initial frame, and let
\(p_{n,m}:X_n^{\mathrm{mar}}\to X_m^{\mathrm{mar}}\) be its
truncations. At step \(j\), the prefix contains a visible
word \(w_j\) and a frame \(f_j\). The frame update
has the form \(f_{j+1}=g_j(x)f_j\), where \(g_j(x)\) depends
only on information already contained in the prefix. A
compatible prolongation therefore preserves all previous
frames.

\begin{proposition}[Naturality of incidence transport]
\label{prop:lf-incidence-naturality}
The map
\begin{equation}
 \mathcal Q_n(x)=
 \bigl((f_j,\gamma_{f_jA_Wf_j^{-1}}(w_j))\bigr)_{0\leq j<n}
 \label{eq:lf-Qn}
\end{equation}
satisfies \(r_{n,m}\mathcal Q_n=\mathcal Q_mp_{n,m}\), where \(r_{n,m}\) truncates
to the first \(m\) encoded components. It therefore induces a
map between the inverse limits of the prefix
and marked-incidence systems.
\end{proposition}
\begin{proof}
Two compatible prefixes have the same words and the same
transport operators up to step \(m-1\). Starting from the
same initial frame, the recursive relation gives by induction the
same \(f_j\) up to that step. The first
\(m\) components of \eqref{eq:lf-Qn} therefore agree. The universal
property of the inverse limit produces the prolongation map.
\end{proof}

In each frame, the first projection of \(\gamma_A\) exactly
recovers the visible word. This scope is local: deep
memory is preserved in the enriched state, and its
recovery requires the corresponding coordinates. Passing
from words to supports retains even less information.
Likewise, a general change of basis in
\(\operatorname{GL}_6(\mathbb F_3)\) also transports the Hamming
metric and the incidence of the chart. Monomial changes
preserve them directly. The supports used in this
section are computed in the explicit chart \eqref{eq:lf-aw}.

The preceding naturality makes explicit the incidence reading that
accompanies the numerical cylinders. Cylinder nesting
and restriction of \eqref{eq:lf-Qn} act on the same
compatible prefix, preserving respectively its positional
evaluation and its marked relations. The ambient space
\(\mathbb F_3^6\), of \(729\) words, provides the linear
domain of the encoding; it remains distinct from
the \(243\) visible words and the \(468\) emissions in the
admissible APP--TRIT--TPK census.

\subsection{Self-dual ternary code and enumeration of its weights}
\label{subsec:lf-code}

\begin{definition}[Code of the dodecaphasic chart]
The linear code associated with \eqref{eq:lf-aw} is
\[
 C_W=\{(w,wA_W):w\in\mathbb F_3^6\}\subset\mathbb F_3^{12}.
\]
The weight \(\operatorname{wt}(c)\) counts the nonzero coordinates
of \(c\), and \(\operatorname{supp}(c)\subset\{1,\ldots,12\}\)
is its set of nonzero positions. A hexad is the support
of a word of weight six; their collection is denoted by \(\mathcal H\).
\end{definition}

\begin{theorem}[Self-duality, minimum distance, and enumerator]
\label{thm:lf-code}
The code \(C_W\) is self-dual with parameters \([12,6,6]_3\) and
has enumerator
\begin{equation}
 W_{C_W}(z)=1+264z^6+440z^9+24z^{12}.
 \label{eq:lf-enumerator}
\end{equation}
\end{theorem}
\begin{proof}
The generator matrix \([I_6\mid A_W]\) has rank six.
By the symmetry of \(A_W\) and \eqref{eq:lf-aw},
\[
 [I_6\mid A_W][I_6\mid A_W]^{\mathsf T}
       =I_6+A_W^2=0.
\]
Thus the code is self-orthogonal of dimension half its
length and is consequently self-dual.

The remaining enumeration is a finite integer operation on
the explicit matrix. For \(r\in\{0,\ldots,6\}\), we group
the words by \(\operatorname{wt}(w)=r\). The following table
gives the number of words of each total weight:
\[
\begin{array}{c|rrrr|r}
r&0&6&9&12&\text{total}\\ \hline
0&1&0&0&0&1\\
1&0&12&0&0&12\\
2&0&60&0&0&60\\
3&0&120&40&0&160\\
4&0&60&180&0&240\\
5&0&12&180&0&192\\
6&0&0&40&24&64\\ \hline
\text{total}&1&264&440&24&729
\end{array}
\]
Each row contains \(\binom6r2^r\) words. To verify
all its entries, let \(q\) be an integer from \(0\) to \(728\), and
define its six ternary digits
\[
 d_i(q)=\left\lfloor\frac{q}{3^{6-i}}\right\rfloor\bmod3,
 \qquad
 e_j(q)=\left(\sum_{i=1}^{6}d_i(q)(A_W)_{ij}\right)\bmod3.
\]
The row and column of the contribution of \(q\) are,
respectively,
\[
 r(q)=\sum_i\mathbf1_{\{d_i(q)\ne0\}},\qquad
 h(q)=r(q)+\sum_j\mathbf1_{\{e_j(q)\ne0\}}.
\]
The accumulation recurrence starts with
\(T^{(0)}_{r,h}=0\) and applies
\[
 T^{(q+1)}_{r,h}
 =T^{(q)}_{r,h}
   +\mathbf1_{\{r=r(q)\}}\mathbf1_{\{h=h(q)\}}
 \quad(0\leq q<729).
\]
It produces exactly the table shown. Ternary expansion
gives a bijection between these \(729\) integers and all words
in \(\mathbb F_3^6\), so the operation exhausts the code.
Summing its rows proves \eqref{eq:lf-enumerator}; the first
nonzero weight is six.
\end{proof}

\begin{remark}[Reproducible evaluation of the finite recurrence]
The following loop literally implements the sums in the proof.
Its only data are the integers of the matrix
\eqref{eq:lf-aw}; it computes the enumerator before comparing it with
the equality obtained.
\end{remark}
{\small
\begin{verbatim}
from collections import Counter
from itertools import product
A = ((0,1,1,1,1,1), (1,0,1,2,2,1),
     (1,1,0,1,2,2), (1,2,1,0,1,2),
     (1,2,2,1,0,1), (1,1,2,2,1,0))
rows = [Counter() for _ in range(7)]
for w in product(range(3), repeat=6):
    v = tuple(sum(w[i]*A[i][j] for i in range(6)) % 3
              for j in range(6))
    r = sum(x != 0 for x in w)
    h = r + sum(x != 0 for x in v)
    rows[r][h] += 1
total = sum(rows, Counter())
if dict(total) != {0:1, 6:264, 9:440, 12:24}:
    raise ArithmeticError("Incompatible enumerator")
\end{verbatim}
}

\begin{proposition}[Hexads and the Witt design]
\label{prop:lf-witt}
The collection \(\mathcal H\) contains \(132\) hexads, and every
five-element set of positions belongs to a unique hexad.
\end{proposition}
\begin{proof}
Two words of weight six with the same support differ by
a scalar. Indeed, among their six ratios of nonzero
coordinates, one of the signs \(1,-1\) occurs at least three
times. Subtracting the corresponding multiple, if the words
were not proportional, would yield a nonzero word
of weight at most three, contradicting
Theorem~\ref{thm:lf-code}. Each support thus corresponds to
the pair \(\{c,-c\}\), and there are \(264/2=132\) supports.

If two distinct supports shared five positions,
one of the two ratios would occur at least three times
in that intersection. The linear combination cancelling those
coordinates would be supported on a union of size seven
and have weight at most four. The minimum distance rules this out.
A five-element set therefore belongs to at most one hexad.
Since
\[
 132\binom65=792=\binom{12}{5},
\]
it belongs to exactly one.
\end{proof}

The constructed object is the extended ternary Golay code
with its small Witt design. Recognition occurs
after encoding, self-duality, and the distance
computation. The lattice construction requires only
the properties already proved and the marked flag
obtained next.

\subsection{Dodecaphasic flag and origin selection}
\label{subsec:lf-flag}

The frame of the generated record preserves the words
\[
 w_\pi=010211,\qquad
 w_e=201101,\qquad
 w_\varphi=121200,\qquad
 w_{\mathrm{APP}}=022020.
\]
The subscripts \(\pi,e,\varphi\) identify the corresponding
reading regions of the constants; here \(w_e\) preserves
the Euler region. Applying \(\gamma_{A_W}\) and taking supports gives
\[
\begin{aligned}
 P_\pi&=\{2,4,5,6,7,8,9,10,11\},\\
 H_e&=\{1,3,4,6,9,10\},\\
 H_\varphi&=\{1,2,3,4,7,9\},\\
 H_{\mathrm{APP}}&=\{2,3,5,10,11,12\}.
\end{aligned}
\]
The first support has weight nine; the remaining three are
hexads. The record
\[
 K=(234,543,140,729,659,824,621,58,914,794,146,601)
\]
selects
\[
 B_K=\{j:K_j\geq729\}=\{4,6,9,10\}=H_e\cap P_\pi.
\]
The threshold \(729\) belongs to the block format of the
declared reading. In this particular record, integer
thresholds from \(660\) to \(729\) produce the same support, so
the support preserves the selected incidence,
while the complete record also preserves the blocks
and the threshold of its reading.

The signed pre-carry record is
\[
\begin{split}
 Z_0={}&(7,297,353,-430,-715,-1199,\\
       &\hspace{9mm}-3,286,-894,-528,380,663).
\end{split}
\]
Its negative support is the hexad
\[
 H_\alpha^-=\{4,5,6,7,9,10\}.
\]
The signs are received from the pre-carry state and are not
extracted from the expansion of the final scalar.

\begin{proposition}[Agreement of signed and incidence selection]
\label{prop:lf-flag}
The hexad \(H_\alpha^-\) is the unique \(H\in\mathcal H\) satisfying
\[
\begin{gathered}
 B_K\subset H\subset P_\pi,\\
 (|H\cap H_e|,|H\cap H_\varphi|,
                   |H\cap H_{\mathrm{APP}}|)=(4,3,2).
\end{gathered}
\]
\end{proposition}
\begin{proof}
Encoding the words of weight six, or applying the preceding
finite recurrence to their supports, gives the four
hexads over \(B_K\). Their complementary pairs are
\[
 \{1,3\},\qquad\{2,11\},\qquad\{5,7\},\qquad\{8,12\}.
\]
Only the second and third are contained in \(P_\pi\).
The hexad of the second has intersections of size three
with \(H_\varphi\) and with \(H_{\mathrm{APP}}\), and therefore
fails the last condition. That of the third has
intersections \(4,3,2\), and agrees with \(H_\alpha^-\).
\end{proof}

\begin{lemma}[Origin of the marked flag]
\label{lem:lf-origin}
The origin selected by incidence is
\[
 o_*=(H_e\cap H_\varphi)
                 \setminus(H_\alpha^-\cup H_{\mathrm{APP}})
     =\{1\}.
\]
Selection commutes with every bijection applied
simultaneously to the supports.
\end{lemma}
\begin{proof}
The intersection \(H_e\cap H_\varphi\) is
\(\{1,3,4,9\}\). The subtracted union contains \(3,4,9\)
and excludes \(1\). Equivariance follows because
bijections commute with union, intersection, and difference.
\end{proof}

The flag distinguishes one component among twelve before the
metric step. The radial vector is determined afterwards,
within a declared positive chamber. Combinatorial selection
and the lattice operation are thereby kept
distinct and connected.

\subsection{Ternary gluing of the twelve components}
\label{subsec:lf-glue}

Let \(A_2=\mathbb Za_1\oplus\mathbb Za_2\), with bilinear
form of matrix
\[
 B_2=\begin{pmatrix}2&-1\\-1&2\end{pmatrix},
 \qquad \rho=a_1+a_2,\qquad \rho^2=2.
\]
For \(x=ua_1+va_2\),
\[
 x^2=2(u^2-uv+v^2).
\]
The identity
\(4(u^2-uv+v^2)=(2u-v)^2+3v^2\) shows
positivity. The determinant is three. The three
classes of \(A_2^*/A_2\) are represented by
\[
 \varpi_0=0,\qquad
 \varpi_1=\frac{2a_1+a_2}{3},\qquad
 \varpi_2=\frac{a_1+2a_2}{3}.
\]
Pairing with \(a_1,a_2\) verifies that these
vectors belong to \(A_2^*\), and
\(2\varpi_1-\varpi_2=a_1\). Hence
\(j\mapsto\varpi_j+A_2\) additively identifies
\(\mathbb F_3\) with the discriminant.

For \(c\in C_W\), we define
\(\phi(c)=(\varpi_{c_1},\ldots,\varpi_{c_{12}})\)
and
\begin{equation}
 N=\bigcup_{c\in C_W}\bigl(A_2^{12}+\phi(c)\bigr).
 \label{eq:lf-glue}
\end{equation}

\begin{theorem}[Even unimodular glue lattice]
\label{thm:lf-niemeier}
The set \(N\) is a positive, even,
unimodular lattice of rank \(24\). Its norm-two vectors
are exactly the roots of \(A_2^{12}\).
\end{theorem}
\begin{proof}
Linearity of the code and additivity of the discriminant
make the union a subgroup of
\((A_2^*)^{12}\) containing \(A_2^{12}\).
It is therefore a lattice of rank \(24\).

The pairing of discriminant classes is
\(2cd/3\) modulo \(\mathbb Z\) in each component.
Self-orthogonality of \(C_W\) makes their sum vanish and proves
integrality of \(N\). The norm of a class
represented by \(\phi(c)\) is
\(2\operatorname{wt}(c)/3\) modulo \(2\mathbb Z\).
All weights in \eqref{eq:lf-enumerator} are
multiples of three, so \(N\) is even.

Its index over \(A_2^{12}\) is \(|C_W|=3^6\).
Using the determinant of the Gram matrix,
\[
 \det N=\frac{\det(A_2^{12})}{[N:A_2^{12}]^2}
       =\frac{3^{12}}{(3^6)^2}=1.
\]
The ambient form is the orthogonal sum of twelve positive
forms; this proves positivity.
Integrality and determinant one imply
\(N^*=N\).

In a nonzero class of the local discriminant, a
vector has the form \((ua_1+va_2)/3\), with residues
\((u,v)\equiv(2,1)\) o \((1,2)\pmod3\).
The integer \(u^2-uv+v^2\) is positive and divisible
by three, so it is at least three. The representatives
\(\varpi_1,\varpi_2\) attain that value: the minimum
norm of each nonzero class is \(2/3\).
A nonzero codeword has at least
six nonzero coordinates. Every vector in a nontrivial
glue class therefore has norm
at least \(6(2/3)=4\).

In \(A_2^{12}\), norm two can only come
from a single component, because each component
has minimum positive even norm two. The solutions
of \(u^2-uv+v^2=1\) are
\((\pm1,0),(0,\pm1),(1,1),(-1,-1)\).
These are the six roots of \(A_2\), and complete the census
of roots of \(N\).
\end{proof}

The preceding properties identify \(N\) as the
Niemeier lattice with root system \(A_2^{12}\).
The description of all its roots is needed to
control which ones survive passage to the neighbour.

\subsection{Radial chamber and rootless neighbour}
\label{subsec:lf-neighbour}

We fix the radial chamber
\[
 v(a)=\sum_{j=1}^{12}a_j\rho_j,\qquad
 a_j=1+3b_j,\qquad b_j\in\mathbb Z_{\geq0},
\]
where \(\rho_j\) is the copy of \(\rho\) in
component \(j\). Evenness of an index-three neighbour
requires \(v(a)^2\equiv0\pmod{18}\).
Since
\[
 v(a)^2=2\sum_j(1+3b_j)^2
       =24+12\sum_jb_j+18\sum_jb_j^2,
\]
the congruence is equivalent to \(\sum_jb_j\equiv1\pmod3\).
The minimum norm within the chamber is attained with
a single \(b_j=1\) and all others zero, and equals \(54\).
The origin \(o_*=1\) of Lemma~\ref{lem:lf-origin}
selects
\begin{equation}
 v_\alpha=4\rho_1+\rho_2+\cdots+\rho_{12},
 \qquad v_\alpha^2=54.
 \label{eq:lf-valpha}
\end{equation}
The qualification ``minimum'' refers to the declared
chamber. In the full residue class,
\((a_1-2a_2,\rho,\ldots,\rho)\) represents the
same class modulo three and has norm
\(14+11\cdot2=36\). This distinction preserves
the condition determining the coefficient four.

We define
\begin{equation}
\begin{split}
 N_0&=\{x\in N:\langle x,v_\alpha\rangle\equiv0\pmod3\},\\
 \Lambda&=N_0+\mathbb Z\,v_\alpha/3.
\end{split}
\label{eq:lf-lambda}
\end{equation}

\begin{theorem}[Even self-dual rootless neighbour]
\label{thm:lf-leech}
The construction \eqref{eq:lf-lambda} is a
positive, even, self-dual lattice of rank \(24\)
with no norm-two vectors. Consequently, it is
isometric to the Leech lattice.
\end{theorem}
\begin{proof}
The character \(x\mapsto\langle x,v_\alpha\rangle\bmod3\)
is nonzero. A simple root of component \(j\)
has pairing \(a_j\equiv1\pmod3\) with
\(a_j\rho_j\). Hence \([N:N_0]=3\), whence
\(\det N_0=9\).

For \(x\in N_0\), the pairing
\(\langle x,v_\alpha/3\rangle\) is an integer, so
\(v_\alpha/3\in N_0^*\). Moreover,
\(\langle v_\alpha,v_\alpha\rangle=54\) implies
\(v_\alpha\in N_0\). If \(v_\alpha/3\in N\),
integrality of \(N\) would force all
pairings of \(N\) with \(v_\alpha\) to be
divisible by three, contradicting the nonzero
character. The class of \(v_\alpha/3\) over \(N_0\)
therefore has order exactly three. It follows that
\[
 [\Lambda:N_0]=3,\qquad
 \det\Lambda=\frac9{3^2}=1.
\]
The pairings of \(N_0\) with the new
generator are integers, and the norm of this generator
is six. The extension is integral. For
\(x\in N_0\) and \(m\in\mathbb Z\),
\[
 \left(x+m\frac{v_\alpha}{3}\right)^2
 =x^2+2m\frac{\langle x,v_\alpha\rangle}{3}
       +6m^2\in2\mathbb Z.
\]
Thus \(\Lambda\) is even; its integrality and
determinant one prove self-duality.
Positivity and rank are inherited from the ambient
Euclidean space.

To exclude roots, we first consider the
class \(N_0\). All roots of \(N\) belong
to \(A_2^{12}\) by Theorem~\ref{thm:lf-niemeier}.
The pairings of its six local roots
with \(a_j\rho_j\) are \(\pm a_j,\pm2a_j\),
congruent to \(\pm1,\pm2\pmod3\).
No root belongs to \(N_0\).

The other two classes of \(\Lambda/N_0\) can
be represented by \(\pm v_\alpha/3\).
In each component, their vectors belong to
one of the classes
\[
 A_2+\varpi_j\pm\rho/3,\qquad j=0,1,2.
\]
The distinguished component has the same
description because \(4\rho/3\equiv\rho/3\pmod{A_2}\).
For the numerators \((u,v)\) in
\((ua_1+va_2)/3\), the six residues are
\[
 (1,1),(0,2),(2,0),(2,2),(1,0),(0,1)\pmod3.
\]
All are nonzero. Positivity of
\(u^2-uv+v^2\) forces it to be at least one,
and the representatives
\((1,1),(0,-1),(-1,0),(-1,-1),(1,0),(0,1)\)
attain that minimum in all six classes.
Each component has norm at least \(2/9\).
A vector in either of the two new
classes has norm at least
\[
 12\cdot\frac29=\frac83>2.
\]
Roots have been excluded in all three classes.
The classification theorem for positive
even unimodular lattices of rank twenty-four
identifies the unique rootless case with the
Leech lattice \cite{ConwaySloane}.
\end{proof}

The final identification uses an external
classification on an object whose properties have already
been proved. Selection of the distinguished
component preserves the flag arising from the
signed record; the subscript \(\alpha\) expresses
that provenance. The norm and coefficients of the
radial vector are obtained from the metric and
the chamber through the stated congruence.

The lattice \(\Lambda\), its bilinear form, its
twelve ambient planes, and the vector \(v_\alpha/3\)
are thereby determined with all the properties
used by the subsequent deformation. The
code also contains the word
\[
 c_*=(1,1,1,1,1,1,2,1,1,1,1,1),
\]
since \(q_*=(1,1,1,1,1,1)\) satisfies
\(q_*A_W=(2,1,1,1,1,1)\). This is the
full-support datum that allows us to construct and
prove stability of the order-three
isometry in Section~\ref{sec:leech-dualidad}.
